%!TeX root=../main.tex
% Figura 2.2: fluxo de desenvolvimento de um patchset.
% Adaptada de Tadokoro, Siqueira e Meirelles (2025), Fig. 1 ("Concept map of
% the patchset development workflow"). Mantém as etapas e os laços de falha;
% traduz os rótulos e destaca as duas etapas que o kw registra (compilação e
% instalação), retomadas na seção 2.4.
\begin{tikzpicture}[
    x=1cm, y=1cm,
    every node/.style={font=\footnotesize, align=center},
    etapa/.style={draw, rounded corners=3pt, minimum width=2.5cm,
                  minimum height=0.95cm, inner sep=2pt, fill=white},
    reg/.style={etapa, very thick, fill=black!12},
    borda/.style={etapa, rounded corners=9pt, fill=black!4},
    seta/.style={-{Latex[length=1.6mm]}, semithick},
    falha/.style={-{Latex[length=1.6mm]}, semithick, dashed},
    rot/.style={font=\scriptsize\itshape, fill=white, inner sep=1pt},
  ]
  % Linha de cima: da concepção à instalação
  \node[borda] (conc) at (0.0, 2.6) {concepção do\\patchset};
  \node[etapa] (impl) at (3.4, 2.6) {implementação};
  \node[reg]   (comp) at (6.8, 2.6) {compilação\\(\textit{build})};
  \node[reg]   (inst) at (10.2, 2.6) {instalação\\(\textit{deploy})};

  % Linha de baixo: dos testes ao merge
  \node[etapa] (test) at (10.2, 0.0) {testes\\(automáticos e manuais)};
  \node[etapa] (envi) at (6.8, 0.0) {envio aos\\mantenedores e listas};
  \node[etapa] (revi) at (3.4, 0.0) {revisão};
  \node[borda] (merg) at (0.0, 0.0) {integração\\à árvore};

  % Caminho de sucesso
  \draw[seta] (conc) -- (impl);
  \draw[seta] (impl) -- (comp);
  \draw[seta] (comp) -- (inst);
  \draw[seta] (inst) -- (test);
  \draw[seta] (test) -- (envi);
  \draw[seta] (envi) -- (revi);
  \draw[seta] (revi) -- (merg);

  % Laços de falha, todos de volta à implementação
  \draw[falha] (comp.north) -- ++(0, 0.5) -| ([xshift=-4mm]impl.north);
  \node[rot, anchor=south] at (5.1, 3.55) {falha na compilação};
  \draw[falha] (inst.north) -- ++(0, 1.1) -| ([xshift=4mm]impl.north);
  \node[rot, anchor=south] at (6.8, 4.15) {falha na instalação};
  \draw[falha] (test.south) |- (1.8, -0.95) |- ([yshift=-2.5mm]impl.west);
  \node[rot, anchor=north] at (6.8, -1.0) {falha nos testes};
  \draw[falha] (revi.north) -- (impl.south)
    node[rot, midway, right=2pt] {alterações pedidas};
\end{tikzpicture}
